\documentclass{book}

% page layout
\usepackage[letterpaper, textwidth=5.0in, textheight=9.0in]{geometry}
\setlength{\oddsidemargin}{3.25in} % making this up
\addtolength{\oddsidemargin}{-0.5\textwidth}
\setlength{\evensidemargin}{\oddsidemargin}
\usepackage{setspace}
\setstretch{1.08}
\renewcommand{\MakeUppercase}[1]{\textsf{#1}} % hackety hack
\sloppy\sloppypar\raggedbottom\frenchspacing

% text macros
\newcommand{\chapref}[1]{\chaptername~\ref{#1}}
\newcommand{\figref}[1]{\figurename~\ref{#1}}

\title{\bfseries Fundamental Astronomy for Physicists}
\author{David W. Hogg}
\date{March 2026}

\begin{document}

\maketitle
\frontmatter
\chapter{Preface}
This book is about the measurement of distant objects with telescopes.
The question is:
What is it, precisely and rigorously, that we can know about the Universe, given only observations with telescopes?
Of course we won't \emph{only} have observations;
we will also have knowledge of basic physics, and especially the relativities (special and general).
We won't have any gravitational waves or neutrinos or other massless (or effectively massless) information carriers.
This is a book about traditional electromagnetic-messenger astrophysics.
That is, what does the study of light, impinging on an observatory, reveal about the cosmos?

The target audience is physicists who want to understand the Universe and its contents, physically.
My expectation is that you know and understand basic special relativity, and the rudiments of quantum mechanics (because, for example, light is both a particle and a wave).
Basic undergraduate physics prepares the reader to read this book.

I was motivated to start writing this book for a number of reasons:
\begin{itemize}
    \item HOGG: Dark matter yada yada.
    \item HOGG: Great precision of the cosmological model...
    \item HOGG: Coordinate systems in stellar dynamics
    \item HOGG: Physicists get it consistently wrong...?
\end{itemize}

\paragraph{Acknowledgements:}
It is a pleasure to thank
  Mike~Blanton (Carnegie),
  Doug~Finkbeiner (Harvard),
  Gerry~Neugebauer (deceased), and
  Ned~Wright (NASA)
for conversations about these matters over the years.

\tableofcontents
\mainmatter
\chapter{What does a telescope do?}
When we first learn geometric optics in school, we learn how a telescope works.
My version of the standard diagram is in \figref{fig:telescope}.
The most important things about this \emph{highly idealized} diagram are the following:
\begin{itemize}
    \item Parallel rays entering the telescope aperture come to a focus at a single point on the focal plane.
    \item Parallel lines coming from two distinct directions come to foci at two distinct points on the focal plane.
\end{itemize}
That is, the telescope maps incoming-photon \emph{angles} to focal-plane \emph{positions}.
This standard diagram (\figref{fig:telescope}) does not capture several other important points about the (idealized) telescope.
For example, it does not give any sense of the following:
\begin{itemize}
    \item The ideal telescope, which never absorbs any photons prior to their arrival at the focal plane, conserves photon phase-space density, or specific intensity, or surface brightness. This is the subject of \chapref{chap:intensity}.
    \item There are critical diffraction effects that limit the angular resolution of the image and create other telescope-aperture effects in the focal plane. This is the subject of \chapref{chap:physical}.
\end{itemize}
The first of these points relates to unitarity in electromagnetism, and will end up giving us a definition of the concept of ``tranparency.''
We'll expand a bit on the latter of these points right now.

HOGG: quantum!

\chapter{Intensity and transparency}\label{chap:intensity}
HOGG: write

\chapter{Fluxes and fluences}

\chapter{Angles}

\chapter{Celestial coordinate systems}

\chapter{Parallaxes and proper motions}

\chapter{Doppler shifts}

\chapter{Foreground and background separations}

\chapter{Instrument calibration}

\chapter{Confusion}

\chapter{Physical optics and interferometry}\label{chap:physical}

\end{document}
